\newpage

\section{РАЗРАБОТКА СИСТЕМЫ ЗАМЕНЫ И МОДИФИКАЦИИ ОБЪЕКТОВ В ВИДЕО}
\subsection{FateZero \texorpdfstring{\cite{fatezero}}{FateZero}}

В прошлой главе мы рассмотрели арзитектуры нейронных сетей для решения данной задач. По итогу я остановлся на модифицированной архитектуре FateZero, в котором я заменю обычную диффузионну модель на диффузионную модель с обним шагом диффузии. Рассмотри для начала оригим оригинальную архитектуру FateZero.

В данной статье представлен новый подход к редактированию видео с нулевым обучением на основе предварительно обученных диффузионных моделей. Метод, названный FateZero, позволяет редактировать реальные видео с помощью текстовых запросов без дополнительного обучения или использования специальных масок. Рассмотрим основные аспекты этого подхода:

\subsubsection{Основная концепция}

FateZero использует предварительно обученные диффузионные модели для генерации изображений по текстовому описанию. Ключевая идея заключается в адаптации этих моделей для редактирования видео с сохранением временной согласованности между кадрами. Это достигается за счет нескольких нововведений:

\begin{enumerate}
    \item Сохранение карт внимания при инверсии видео
    \item Слияние карт внимания при генерации отредактированного видео
    \item Реформирование механизма самовнимания для учета пространственно-временных зависимостей
\end{enumerate}


\subsubsection{Архитектура и процесс работы}

\paragraph{Инверсия видео}

Процесс начинается с инверсии исходного видео с помощью детерминированной инверсии DDIM (Denoising Diffusion Implicit Models). На этом этапе видео преобразуется в шумовое латентное представление:

\begin{enumerate}
    \item Видео кодируется в латентное пространство с помощью энкодера.
    \item Применяется инверсия DDIM для получения шумового латентного представления.
    \item На каждом шаге инверсии сохраняются карты пространственно-временного самовнимания и кросс-внимания.
\end{enumerate}

\paragraph{Редактирование}

После инверсии происходит редактирование видео на основе целевого текстового запроса:

\begin{enumerate}
    \item Шумовое латентное представление постепенно очищается от шума с учетом целевого запроса.
    \item На каждом шаге очистки от шума происходит слияние текущих карт внимания с сохраненными картами из процесса инверсии.
    \item Для слияния используется специально разработанный блок Attention Blending Block.
\end{enumerate}

\paragraph{Блок слияния внимания (Attention Blending Block)}

Этот блок играет ключевую роль в сохранении структуры и движения исходного видео:

\begin{enumerate}
    \item  Карты кросс-внимания для неизменяемых слов в запросе заменяются на исходные.
    \item  Карты самовнимания сливаются с использованием адаптивной пространственной маски.
    \item  Маска получается из карт кросс-внимания исходного запроса и определяет области для редактирования.
\end{enumerate}

\paragraph{Пространственно-временное внимание}

Для обеспечения согласованности между кадрами авторы модифицируют механизм самовнимания:

\begin{enumerate}
    \item  Стандартные блоки самовнимания заменяются на пространственно-временные.
    \item  Это позволяет учитывать зависимости не только в пределах кадра, но и между кадрами.
\end{enumerate}

\subsubsection{Преимущества подхода}

\begin{enumerate}
    \item  Нулевое обучение: Метод не требует дополнительного обучения для каждого нового запроса на редактирование.
    \item  Отсутствие пользовательских масок: Нет необходимости в ручной разметке областей для редактирования.
    \item  Сохранение структуры и движения: Благодаря слиянию карт внимания сохраняется исходная структура видео.
    \item  Универсальность: Подход работает как с моделями для генерации изображений, так и с видео-моделями.
\end{enumerate}

\subsubsection{Применения}

FateZero демонстрирует возможности в различных задачах редактирования видео:

\begin{enumerate}
    \item  Изменение стиля видео (например, преобразование в акварельную анимацию)
    \item  Локальное редактирование атрибутов объектов
    \item  Замена объектов с сохранением формы и движения
    \item  Редактирование отдельных изображений
\end{enumerate}


\subsection{Stable Diffusion XS \texorpdfstring{\cite{sdxs}}{Stable Diffusion XS}}

Для ускорения FateZero вместо Stable Diffusion 1.4 из оригинальной статьи будет исопозоваться подход по оптимизации архитектуры модели из Stable Diffusion XS. Рассмотрим методы оптимизации и дистиляции SDXS.

\subsubsection{Архитектурные оптимизации}

\paragraph{Оптимизация декодера изображений}

Авторы начинают с оптимизации декодера изображений, который является частью вариационного автоэнкодера (VAE) в архитектуре Latent Diffusion Model (LDM). Они предлагают метод дистилляции VAE (VAE Distillation, VD), который обучает легковесный декодер изображений для имитации выхода оригинального декодера VAE. Функция потерь для этого процесса включает в себя:


\begin{enumerate}
    \item  Потери реконструкции: L1-норма между сгенерированным и целевым изображением, измеренная на 8-кратно уменьшенных версиях изображений.
    \item  Потери GAN: для улучшения качества и реалистичности генерируемых изображений.
\end{enumerate}

Функция потерь VD выглядит следующим образом:


где G - это обучаемый легковесный декодер, z - латентный код, ˜x - изображение, реконструированное оригинальным декодером VAE, а D - дискриминатор GAN.

\paragraph{Оптимизация U-Net}

Далее авторы обращаются к оптимизации основной части диффузионной модели - архитектуры U-Net. Они применяют стратегию удаления блоков, вдохновленную работой BK-SDM. Этот подход включает в себя:

\begin{enumerate}
    \item  Селективное удаление резидуальных и трансформерных блоков из U-Net.
    \item  Обучение более компактной модели для воспроизведения промежуточных карт признаков и выходных данных оригинальной модели.
\end{enumerate}

Процесс дистилляции знаний для U-Net использует две функции потерь:



\paragraph{Оптимизация ControlNet}

Авторы также предлагают метод дистилляции для ControlNet, который является расширением U-Net для задач генерации изображений с условиями. Процесс дистилляции ControlNet включает:

\begin{enumerate}
    \item Копирование архитектуры энкодера U-Net с добавлением дополнительных сверточных слоев для включения пространственного управления.
    \item Дистилляцию ControlNet оригинального U-Net в соответствующий ControlNet уменьшенного U-Net.
    \item Комбинирование ControlNet с U-Net для дистилляции промежуточных карт признаков и выходных данных U-Net.
\end{enumerate}

\subsubsection{Стратегия обучения одношаговой модели}

Ключевой инновацией предложенного подхода является метод обучения модели для генерации изображений за один шаг. Авторы разработали быструю и стабильную стратегию обучения, которая включает в себя:

\begin{enumerate}
    \item Выпрямление траектории сэмплирования.
    \item Быстрая доводка многошаговой модели до одношаговой путем замены функции потерь дистилляции на предложенную функцию потерь сопоставления признаков.
    \item Расширение стратегии обучения Diff-Instruct, используя градиент предложенной функции потерь сопоставления признаков для замены градиента, предоставляемого дистилляцией оценки во второй половине временного шага.
\end{enumerate}

Этот подход основан на концепции интегрального расхождения Кульбака-Лейблера (IKL) между двумя распределениями:

 обозначает образец, сгенерированный обучаемым одношаговым генератором, а s и spt обозначают функции оценки DM, обученной онлайн на сгенерированных данных, и предварительно обученной DM соответственно.
